(1) 均匀细梁的材料杨氏模量 $E$、截面二次矩 $I$、质量密度 $\rho$、截面积 $A$，忽略转动惯性与剪切形变，小横位移满足 $\rho A y_{tt}+EI y_{xxxx}=0$。求无限梁的实波数 $k$ 色散关系。
(2) 计算弯曲波的相速度与群速度，并比较其频率依赖关系。
(3) 长 $L$ 的两端简支梁满足 $y(0,t)=y(L,t)=0$ 且弯矩零 $y_{xx}(0,t)=y_{xx}(L,t)=0$，求本征模与离散频率。
(4) 写出梁的波动总能量并证明简支边界下能量守恒，检验加长梁对最低频率的影响。
